% Generated draft: requires mathematical and source review.
\documentclass[11pt]{article}
\usepackage[margin=25mm]{geometry}
\usepackage{fontspec}
\setmainfont{Latin Modern Roman}
\usepackage{amsmath,amssymb,mathtools,unicode-math}
\setmathfont{Latin Modern Math}
\usepackage{enumitem,needspace,xurl,hyperref}
\setlength{\parindent}{0pt}
\setlength{\parskip}{5pt}
\setlength{\emergencystretch}{3em}
\newcommand{\N}{\mathbb{N}}
\newcommand{\Z}{\mathbb{Z}}
\newcommand{\Q}{\mathbb{Q}}
\newcommand{\R}{\mathbb{R}}
\newcommand{\C}{\mathbb{C}}
\newcommand{\F}{\mathbb{F}}
\newcommand{\A}{\mathbb{A}}
\newcommand{\PP}{\mathbb{P}}
\newcommand{\E}{\mathbb{E}}
\newcommand{\eps}{\varepsilon}
\newcommand{\dd}{\,\mathrm{d}}
\begin{document}
\section*{UnsolvedMath ID: 700665}
\textbf{Source problem code:} catalogue entry
\subsection{Definitions and mathematical statement}
Let \(G\) be a finite group, let \(\operatorname{Irr}_{\C}(G)\) denote its irreducible complex characters, and let \(1_G\) denote its principal character, so that \(1_G(g)=1\) for every \(g\in G\). A linear character of a cyclic subgroup \(C\leq G\) is a homomorphism \(\lambda:C\to\C^{\times}\). Its induced character is
\[
\bigl(\operatorname{Ind}_{C}^{G}\lambda\bigr)(g)
=\frac{1}{|C|}\sum_{\substack{x\in G\\x^{-1}gx\in C}}\lambda(x^{-1}gx),
\qquad g\in G.
\]
These are the \emph{induced cyclic characters} of \(G\). The cyclic subgroups and their linear characters are unrestricted; in particular, the trivial subgroup and principal linear characters are included.

Describe all finite groups \(G\) with the following property: for every \(\chi\in\operatorname{Irr}_{\C}(G)\), there exist integers \(a_\chi\) and \(b_{\chi,C,\lambda}\) such that, as class functions on \(G\),
\[
\chi=a_\chi 1_G+
\sum_{\substack{C\leq G\\C\text{ cyclic}}}
\ \sum_{\lambda\in\operatorname{Hom}(C,\C^{\times})}
 b_{\chi,C,\lambda}\operatorname{Ind}_{C}^{G}\lambda.
\]
The coefficients may depend on \(\chi\) and may be negative; both sums are finite.

The supplied source reports that computer-aided computations of A. V. Rukolaine establish this property for \(\mathbb{S}_5\), \(\mathrm{SL}_2(13)\), and \(M_{11}\), respectively the symmetric group on five letters, the determinant-one matrix group over \(\F_{13}\), and the Mathieu group \(M_{11}\). This computational assertion is retained as source-reported background and has not been independently verified here.
\subsection{Short English statement}
Classify the finite groups whose every irreducible complex character is an integer linear combination of the principal character and characters induced from linear characters of cyclic subgroups. Negative coefficients are allowed. The supplied registry labels this problem open; this draft does not independently verify that status.
\subsection{Sources}
\begin{itemize}
\item Local UnsolvedMath registry, supplied complete record with permanent ID \(700665\), exact source code \texttt{catalogue entry}, and title \emph{Kourovka Notebook Problem 10.50}; supplied update date: 2026-10-16. The registry status is \texttt{open}, a supplied label rather than an independently verified literature assessment.
\item \emph{The Kourovka Notebook (Unsolved Problems in Group Theory)}, original problem 10.50, 10th issue (1986), proposed by A. V. Rukolaine. The supplied background identifies the 21st-issue edition (2026), arXiv:1401.0300v48; the collection description names E. I. Khukhro and V. D. Mazurov as editors.
\item A. V. Rukolaine, \emph{Abstracts of the 10th All--Union Symposium on Group Theory}, Gomel', 1986, p.~199 (Russian), as cited in the supplied problem statement for the computer-aided computations and examples. This reference and the Notebook text have not been independently consulted here.
\end{itemize}
\end{document}
